\documentclass[11pt]{article}
\usepackage[margin=1in]{geometry}
\usepackage{amsmath,amssymb,amsthm}
\usepackage{booktabs}
\usepackage{hyperref}

\title{A Formalization Plan: Machine-Checkable Verification of the Schwarzschild Solution\\[4pt]
\large Stage 5 pilot of the Computational Attack Program on Fundamental Unification}
\author{Kimi Work research scaffold}
\date{2026-09-16}

\begin{document}
\maketitle

\section{Why this target}
A theory-of-everything program lives or dies on whether its claims can be \emph{checked},
not asserted. Before formalizing anything exotic, the pipeline must demonstrate it can
formalize something \emph{known}. The Schwarzschild verification is the ideal pilot:
it is a finite computation, its correctness is undisputed, and every step is
mechanizable.

\begin{quote}
\textbf{Claim to be formalized.} The metric
\[
  ds^2 = -\left(1-\frac{2GM}{r}\right)dt^2
        + \left(1-\frac{2GM}{r}\right)^{-1}dr^2
        + r^2\left(d\theta^2 + \sin^2\theta\,d\phi^2\right)
\]
satisfies the vacuum Einstein equations $R_{\mu\nu} = 0$ for $r > 2GM$.
\end{quote}

\section{Formalization ladder}
\begin{enumerate}
  \item \textbf{Level 0 --- symbolic check (sympy).} Compute Christoffel symbols
        $\Gamma^{\rho}_{\mu\nu}$, the Ricci tensor $R_{\mu\nu}$, and verify
        $R_{\mu\nu} \equiv 0$ componentwise. Fully automatable today.
  \item \textbf{Level 1 --- scripted differential geometry.} Re-derive the same result
        from an index-free formulation (orthonormal coframe, Cartan structure equations
        $d\theta^a + \omega^a{}_b \wedge \theta^b = 0$, curvature 2-forms
        $\Omega^a{}_b = d\omega^a{}_b + \omega^a{}_c \wedge \omega^c{}_b$).
  \item \textbf{Level 2 --- proof assistant (Lean 4 + Mathlib).} Formalize the claim
        against Mathlib's manifold and Riemannian geometry infrastructure. The theorem
        statement must define the metric, the Ricci tensor, and the domain
        $r > 2GM$ --- the proof itself is then a computation the kernel can check.
\end{enumerate}
Each level must reproduce the result of the level below it; a disagreement at any
level halts the pipeline and is itself a reportable finding.

\section{Deliverables and acceptance criteria}
\begin{center}
\begin{tabular}{@{}lll@{}}
\toprule
Deliverable & Form & Acceptance criterion \\
\midrule
D1: symbolic verification & Python script & $R_{\mu\nu}=0$ verified symbolically, all 10 components \\
D2: Cartan re-derivation & LaTeX note & compiles; every step annotated \\
D3: Lean formalization & Lean 4 project & \texttt{lake build} succeeds, zero \texttt{sorry} \\
D4: public artifact & GitHub repo & third party can rerun D1--D3 \\
\bottomrule
\end{tabular}
\end{center}

\section{Why this matters for the larger program}
Recent agentic formalization work --- e.g. \emph{Hilbert} (ICLR 2026), the
\emph{Ax-prover} framework for mathematics and quantum physics (2026), and
\emph{FormaTheoria}'s progress toward the classification of finite simple groups ---
shows that theory-scale Lean formalization is now feasible. If this pilot succeeds,
the same ladder is applied to the program's real targets: the Einstein--Hilbert
functional-RL truncation used in asymptotic safety, and the Benincasa--Dowker curvature
operator on causal sets. If the pilot fails, the failure localizes exactly which
mathematical infrastructure is missing --- also a useful, publishable negative result.

\section{Kill criteria}
The pilot is abandoned if: (a) Level 0 and Level 1 disagree; (b) Level 2 requires
axioms beyond standard Mathlib without justification; or (c) no component of the
verification completes within the allotted effort --- in which case the bottleneck,
not the theorem, becomes the object of study.

\end{document}
